\section{Projection for naturally longer output}
\label{app:forced}

The main budget table leaves AR near its natural EOS at 126 tokens
while charging DG for a specified number of full canvases. A different
question is whether a task that requires multiple records would make both
models generate longer outputs. Table~\ref{tab:forced} evaluates the
fitted equations at assumed output lengths for that scenario.

\begin{table}[ht]
\centering
\caption{Illustrative fixed-prompt cost-model extrapolation. Only the
single-record anchor is observed (AR mean versus DG median anchor).
Multi-record lengths and costs are assumed/calculated, not measurements
of multi-record extraction. Quality is not evaluated for those rows.}
\label{tab:forced}
\begin{tabular}{@{}lrrrr@{}}
\toprule
Scenario & Assumed $L$ & AR (ms) & DG (ms) & AR/DG \\
\midrule
Single-record anchor & 126 & 945 & 379 & $2.5\times$ \\
Two records & 280 & 2051 & 498 & $4.1\times$ \\
Four records & 560 & 4067 & 618 & $6.6\times$ \\
Eight records & 1120 & 8099 & 857 & $9.5\times$ \\
\bottomrule
\end{tabular}
\end{table}

The calculation assumes about 140 tokens per record and leaves the
reference prompt cost fixed. A real multi-record task also lengthens the
input and may change formatting, stopping behavior, and accuracy.
Consequently these ratios are hypotheses for a new experiment, not
validated speedups or evidence that the observed serial ratio is a
conservative lower bound.
